% ═══════════════════════════════════════════════════════════════════════════
% QR-CODE: Stundenleistung Die Federspannkraft (UE 4)
% Physik Klasse 7 – Beamer-Projektion
% ═══════════════════════════════════════════════════════════════════════════

\documentclass[a4paper,11pt]{article}

\usepackage[utf8]{inputenc}
\usepackage[T1]{fontenc}
\usepackage{helvet}
\renewcommand{\familydefault}{\sfdefault}

\usepackage[margin=20mm]{geometry}
\usepackage{xcolor}
\usepackage{qrcode}
\usepackage{hyperref}
\hypersetup{colorlinks=true, urlcolor=physik}

% Farben
\definecolor{physik}{HTML}{2c5aa0}

% Konfiguration
\newcommand{\mttitel}{Die Federspannkraft}
\newcommand{\mturl}{https://devbox42.github.io/sciencesim/physik/kl07/02-kraefte/04-federspannkraft/MT-04-federspannkraft.html}
\newcommand{\mturlkurz}{devbox42.github.io/.../04-federspannkraft/MT-04-federspannkraft.html}

\pagestyle{empty}

\begin{document}

\begin{center}

% Titel
\vspace*{10mm}
{\fontsize{28}{34}\selectfont\bfseries\color{physik}Stundenleistung: \mttitel}

\vspace{15mm}

% QR-Code mit Rahmen
\fcolorbox{physik}{white}{\qrcode[height=100mm]{\mturl}}

\vspace{12mm}

% URL (klickbar)
{\large\href{\mturl}{\mturlkurz}}

\vspace{8mm}

% Hinweis
{\Large\color{gray}Scanne den QR-Code mit deinem Gerät}

\vspace{10mm}

% Info-Box
\fcolorbox{physik}{physik!10}{%
    \parbox{0.7\textwidth}{%
        \centering
        \vspace{3mm}
        {\large 15 Minuten | Niveau wählen | Auswerten}\\[2mm]
        Wähle dein Niveau und bearbeite die Aufgaben.
        \vspace{3mm}
    }%
}

\end{center}

\end{document}
